\documentclass[11pt]{scrartcl}
\usepackage[sexy]{../evan}
\usepackage{listings}    
\usepackage{xcolor}    

\lstset{
  language=C++,
  basicstyle=\ttfamily\small,
  keywordstyle=\color{blue}\bfseries,
  stringstyle=\color{red},
  commentstyle=\color{green!50!black},
  % The following lines override evan.sty defaults:
  frame=none,           % Removes the shadowbox/border
  numbers=none,         % Removes the line numbers
  backgroundcolor={},   % Makes the background transparent (so theorem color shows)
  xleftmargin=0pt,      % Removes extra indentation
  xrightmargin=0pt
}

\declaretheorem[style=thmbluebox,name=Solution,numberwithin=section]{solution}
\declaretheorem[style=thmredbox,name=Problem,sibling=theorem]{prob}

\begin{document}
\title{CodeForces Review}
\date{3 January 2026}
\maketitle

\begin{center}
  \itshape
  Sometimes figuring out the langauge semantics is harder then implementing the algorithm.
\end{center}

An important skill for olympiad geometers is to be able to guess when three points are collinear,
four points are concyclic, three lines are concurrent, and so on.
Difficult geometry problems (on the level of IMO 3/6) often amount to finding two or three critical claims;
each of these claims may be no harder to prove than an IMO 1/4, but making the right guesses of \emph{what} to prove
can turn out to the core difficulty of the problem.
For a fantastic example, see my solution to
\href{http://www.aops.com/community/c6h418983p3518149}{IMO 2011/6}.

In this exercise, I'll put write down a configuration of several points, lines, and circles.
Your job is to find as many ``coincidences'' as you can: nontrivial collinearities, equal angles,
cyclic quadrilaterals, and so on and so forth.
Then, see if you can prove them!

\tableofcontents

\newpage

\section{900 rated}
\begin{problem}
   B. Chemistry \\

   Given a string $s$ of length $n$, consisting of lowercase latin letters, and an integer $k$. 
   Check if it is possible to remove \textbf{extactly k} characters from the string $s$ in such a way that the remaining 
   can be rearranged to form a palindrome. Note that you can reorder the remaining characters in any way. \\

   A \textbf{palindrome} is a string that is equal to its reversal.
\end{problem}

\begin{solution}
    We only need to check how many letters have an odd occurence in the string. If we have $x$ letters with odd occurences 
    then this cannot be greater than $k + 1$. We are allowed to have 1 occurence of a odd letter in a palindrome or 0. Since we 
    have up to $k$ removables then after $x - k$ we must have odd letters $\leq 1$. 
\begin{lstlisting}
int main() {
  int t;
  cin >> t;
  while (t--) {
    int n, k;
    string s;
    cin >> s;
    vector<int> freq(26);
    for (char c : s) {
      freq[c-'a']++;
    }
    int odd = 0;
    for (int x : freq) {
      if ((x&1)) odd++;
    }
    if (odd > k + 1) {
      cout << "NO\n";
    } else {
      cout << "YES\n";
    }
  }
}
\end{lstlisting}
\end{solution}

\newpage

  \begin{problem}
    A. Two Permutations \\ 

    Given three integers, $n, a,$ and $b$. Determine if there exists two permutations $p$ and $q$ of length $n$, which 
    meet the following conditions: 
      \begin{itemize}
        \item The length of the longest common prefix of $p$ and $q$ is a.
        \item The length of the longest common suffix of $p$ and $q$ is $b$.
      \end{itemize}
      A permutation of length of $n$ is an array containing each integer from $1$ to $n$ exactly once. For 
      example, $[2,3,1,5,4]$, is a permutation. \\ 

      For example for $n = 4, a = 1, b = 1$ we could have $[1,2,3,4]$ or $[1,3,2,4]$.
\end{problem}

\begin{solution}
  If $a + b + 2 \leq n$, then there is always a way to find a valid paring for exmaple:
    \[
        A = \{\textcolor{red}{1,2,\cdots, a,} \textcolor{orange}{n-b,} a+1, a+2, \cdots, n-b-1, n-b+1, n-b+2, \cdots, n\}
    \]
    \[
      B = \{1,2,\cdots,a,a+1,a+2,\cdots,n-b-1,n-b,n-b+1,n-b+2,\cdots, n\}
    \]

    We can seprate the array into the left part, middle, right part, where left is $a$ and rigth is $b$ then the middle there is the
    case the must have atlest two numbers to be able to not grow the prefix or suffix for $A$  and $B$ which is why we have $a + b + 2$ and this 
    cannot be greater than the length of the array. The case where we have $a == b == n$ since we have the original prefrix and suffix equal to the array 
    we are guranteed the longest. 

\begin{lstlisting}
int main() {
  int t;
  cin >> t;
  while (t--) {
    int n, a, b;
    cin >> n >> a >> b;
    if (a+b+2 <= n || (a==b && b == n)) {
      cout << "YES\n";
    } else {
      cout << "NO\n";
    }
  }
}
\end{lstlisting}
\end{solution}

\end{document}